\begin{abstract}
Multi-task model merging has emerged as a powerful paradigm for combining multiple task-specific expert neural networks into a single unified model without the high computational costs of retraining. 
Recent efforts have turned to test-time adaptive merging (e.g., AdaMerging), which optimizes merging coefficients on a small, unlabeled calibration stream at test time. 
However, a fundamental, unstudied question persists: \emph{What is the optimal structural granularity at which these merging coefficients should be defined and optimized, and how does this choices affect multi-task generalization?} 
To answer this, we propose \textbf{GranMerge}, a rigorous empirical framework to deconstruct the \emph{Generalization-Granularity Trade-off} in adaptive model merging. We systematically investigate five nested levels of parameter resolution (Global, Layer-wise, Block-wise, Component-wise, and Tensor-wise) across four visual classification tasks (MNIST, FashionMNIST, CIFAR-10, SVHN) in a resource-constrained, low-fidelity expert regime using a 12-layer Vision Transformer. 
Our exhaustive sweeps across multiple seeds and two distinct optimization families---unconstrained first-order Adam gradient descent and zero-order 1+1 Evolution Strategies (ES)---unveil critical findings:
(1) Increasing structural granularity leads to severe \emph{transductive overfitting} on the compact local test-time calibration batch ($N=256$). Rather than finding superior routing, fine-grained high-dimensional optimization exploits local entropy, degrading global multi-task generalization. 
(2) Optimization dynamics dictate overfitting susceptibility: Adam gradient descent is highly vulnerable to rapid generalization collapse, whereas black-box zero-order 1+1 ES maintains higher generalization, which we empirically deconstruct as a combination of isotropic search boundaries and optimization sluggishness (underfitting) in high-dimensional spaces. 
(3) Simple L2 spatial-depth regularizers (Elastic Spatial Regularization and Total Variation) successfully stabilize the zero-order ES optimizer (improving tensor-wise performance from 29.43\% to 30.17\% for 1+1 ES, and from 26.91\% to 28.51\% for Adam), but are insufficient to arrest the chaotic, unconstrained updates of first-order Adam. 
Our work delivers the first comprehensive empirical mapping of the structural boundaries of adaptive weight blending, providing honest, actionable guidelines for stable multi-task deployment under test-time constraints.
\end{abstract}
